\section{Introduction}
Organizations convey ethical norms in two ways. One is written down in codes of conduct and training modules. The other is practiced: it is visible in which shortcuts are tolerated, which objections are taken seriously, and what happens to people who bend a rule or refuse to. Newcomers learn the first set by reading. They learn the second mainly by watching. A study of new accountants found that newcomers obtained normative information chiefly through observation and chiefly from peers~\citep{morrison1993newcomer}, and research on ethical leadership rests on the premise that followers learn what is expected by observing leaders whose conduct serves as a model~\citep{brown2005ethical}. How well newcomers learn the practiced norms therefore depends on how much of that practice they can see.

For a large group of recent entrants, much less of it was visible. Working from home spread rapidly in 2020~\citep{kniffin2021covid}. When one large firm moved to fully remote work, its employees' collaboration networks became more siloed and their communication more asynchronous~\citep{yang2022effects}, and newcomers onboarded remotely report fewer interactions with insiders~\citep{paik2026membership}. Among those starting their first office-based jobs in this period were many members of Generation Z, as that cohort is commonly defined. It is plausible, although we have found no data that establish it, that a substantial share of them spent part of their entry period at a distance from the colleagues whose conduct would otherwise have taught them the practiced norms.

As far as we can tell, the question has not been asked, in part because the relevant literatures have developed separately. Research on ethical leadership and socialization explains how ethical norms pass from experienced members to newcomers, but it has studied settings in which the two share space~\citep{brown2005ethical}. Research on remote work and remote onboarding documents the loss of everyday interaction, but it examines adjustment, relationships, and performance rather than ethics~\citep{paik2026membership}. Imprinting research shows that conditions at the start of a career leave durable marks, but its outcomes have been performance and professional orientation~\citep{tilcsik2014imprint}. Research on generations, for its part, has concentrated on whether cohorts differ in their values. None of these literatures asks what newcomers learn about workplace ethics when they enter at a distance, or whether it lasts.

This framing also departs from the usual way of connecting Generation Z and workplace ethics, which asks whether the cohort holds different ethical values. The evidence for generational differences in work attitudes and values is weak and difficult to separate from age and period effects~\citep{costanza2012generational,rudolph2021generations}, although the debate is not settled~\citep{campbell2015generational}. We ask instead how members of the cohort learned the ethics of their workplaces. That question returns to Mannheim's view that a cohort becomes a generation through shared formative experiences~\citep{mannheim1952problem}, and to the distinction between generations defined by age and generations defined by entering an organization at the same time~\citep{joshi2010unpacking}.

Two terms need definition. By workplace ethics we mean the standards of conduct that govern how members of an organization treat one another, clients, and others affected by their work, as those standards are both stated and practiced. Research on ethical infrastructure separates the formal elements that state such standards from the informal elements that convey how they are actually applied~\citep{tenbrunsel2003building}, and our focus is on the latter. By an ethically charged situation we mean one in which a decision affects the interests of others and the right course is contested or unclear, such as a request to shade a report or to overlook a colleague's shortcut. By Generation Z we mean the cohort that followed the Millennials into the labor market. Its birth-year boundaries are set by convention and vary across sources, so we use the label descriptively.

We develop a conceptual model of ethical socialization under distributed entry grounded in two established theories. Social learning theory explains how enacted norms are acquired through observation of others and of the consequences they experience~\citep{bandura1977social}. Imprinting theory explains why characteristics formed during a sensitive period persist after conditions change~\citep{marquis2013imprinting}. The model proposes that distributed entry reduces newcomers' clarity about enacted ethical norms by limiting their observational access to ethical conduct, and that leaders who explicitly narrate the ethical reasoning behind their decisions can offset part of this loss. It also proposes a two-sided effect: because distance weakens the transmission of enacted norms in general, it deprives newcomers of good examples in ethical units but also shields them in part from misconduct that has become normal elsewhere. The final propositions concern whether these effects persist after the entry period and whether they follow the conditions of entry rather than birth cohort.

The two-sided prediction is the part of the model with the most practical weight. Discussions of remote work and ethics tend to ask whether distance makes employees behave better or worse. If the model is right, that question has no single answer, because the effect of distance on a newcomer depends on the norms of the unit the newcomer joins. The prediction also distinguishes our account from a simpler one in which distance reduces supervision and so creates more room for misconduct everywhere. The two accounts make different predictions for units in which misconduct is already normal, which gives research a way to tell them apart.

The paper makes three contributions, all of which remain to be tested.

\begin{itemize}
\item It treats ethical socialization as a transmission process whose fidelity depends on observational access, and it introduces enacted-norm clarity and leader ethical narration as constructs for studying that process when newcomers and their models do not share space. Existing ethical leadership research explains role modeling well but assumes co-presence~\citep{brown2005ethical}.
\item It proposes that distance weakens the transmission of unethical as well as ethical norms. Research on the normalization of corruption shows that newcomers can be socialized into misconduct~\citep{ashforth2003normalization}; our model implies that the channel through which this happens can narrow, and it yields a prediction that distinguishes a learning account of remote-work misconduct from an account based on reduced monitoring.
\item It extends imprinting theory, which has linked entry conditions to later performance~\citep{tilcsik2014imprint}, to ethical norms, and it relocates the Generation Z question in the conditions of organizational entry, stated as a proposition that research can confirm or reject.
\end{itemize}
